% Included by slide.tex. Reusable task prompt: AGENT_BRIEF.md.
\begin{frame}[t,plain]
    \titlepage
\end{frame}

\begin{frame}{What a successful contribution contains}
    \begin{itemize}
        \item A module that fits existing discovery and execution.
        \item Correct MISP inputs, outputs, configuration and provenance.
        \item Tests that demonstrate behavior, not just Python syntax.
        \item Useful metadata, documentation and a focused pull request.
    \end{itemize}
    \begin{block}{The agent assists development}
        It is not a new module family or a reason to add a model service
        to the runtime implementation.
    \end{block}
    \note{This companion assumes the MISP modules introduction.
    It is provider-neutral and works with the coding assistant the contributor
    already uses. The reference exercise requires no model credentials or agent SDK.}
\end{frame}

\begin{frame}{Keep ownership and evidence explicit}
    \small
    \begin{tabular}{@{}P{2.2cm}P{4cm}P{4cm}@{}}
        \toprule
        Contributor & Coding agent & Maintainer \\
        \midrule
        Defines the need and boundaries. &
        Inspects code and prepares a plan and patch. &
        Assesses fit and maintenance cost. \\
        Checks contracts and reviews results. &
        Runs permitted checks and reports actual outcomes. &
        Reviews the patch and evidence. \\
        Owns the submitted contribution. &
        Identifies assumptions and limits. &
        Decides whether to merge. \\
        \bottomrule
    \end{tabular}
    \medskip
    A confident explanation from an agent is not execution evidence.
    \note{Agree on allowed network access and writes in the task brief.
    Do not claim a provider was tested when only mocked responses ran.
    Accountability remains with the contributor and maintainers.}
\end{frame}

\begin{frame}{A repeatable development loop}
    \centering
    \begin{tikzpicture}[
        agentstep/.style={draw=mispTeal,fill=mispBlue!5!white,rounded corners,
            align=center,text width=2.7cm,minimum height=1.3cm,font=\small},
        flow/.style={-{Stealth},thick,draw=mispTeal}]
        \node[agentstep] (brief) at (0,2) {1. Specify the gap\\and contract};
        \node[agentstep] (inspect) at (3.7,2) {2. Read code,\\docs and tests};
        \node[agentstep] (patch) at (7.4,2) {3. Implement\\a small patch};
        \node[agentstep] (test) at (7.4,-.4) {4. Run unit and\\service checks};
        \node[agentstep] (review) at (3.7,-.4) {5. Review output,\\diff and limits};
        \node[agentstep] (pr) at (0,-.4) {6. Write docs\\and open a PR};
        \draw[flow] (brief) -- (inspect);
        \draw[flow] (inspect) -- (patch);
        \draw[flow] (patch) -- (test);
        \draw[flow] (test) -- (review);
        \draw[flow] (review) -- (pr);
    \end{tikzpicture}
    \sources{Reference snapshot: misp-modules 3a903305cf37, reviewed 29 October 2026.}
    \note{Repeat the relevant loop when a check reveals a defect.
    Read the current target checkout: a training snapshot or agent memory is not
    the authority for a later API. The README pins the reviewed evidence.}
\end{frame}

\section{Read the existing contract}

\begin{frame}{Give the agent the right repository context}
    \begin{itemize}
        \item Contributor instructions, dependency declarations and CI.
        \item Loader and dispatcher: \code{misp_modules/__init__.py}
        and \code{__main__.py}.
        \item A recent module of the \textbf{same family}, plus its tests.
        \item Existing shared libraries and target-format builders.
        \item Official API/schema documentation, version and examples.
    \end{itemize}
    \begin{block}{Is a new module needed?}
        Search existing capabilities before implementing a duplicate integration.
    \end{block}
    \note{The agent should report the checkout revision and reference files.
    A marketing description is not a wire contract. If official details are missing,
    record the ambiguity and continue independent work rather than inventing behavior.}
\end{frame}

\begin{frame}{Select a family before writing code}
    \small
    \begin{tabular}{@{}P{2.5cm}P{2.3cm}P{5.4cm}@{}}
        \toprule
        Directory & Metadata type & Typical contract \\
        \midrule
        \code{expansion} & expansion / hover & Observable or attribute to enrichment results. \\
        \code{import_mod} & import & File, paste or report to parsed MISP data. \\
        \code{export_mod} & export & Supplied MISP data to an encoded artifact. \\
        \code{action_mod} & action & Workflow data and parameters to an operation/result. \\
        \bottomrule
    \end{tabular}
    \medskip
    Match the specific module's schema; this is not one universal request format.
    \note{The filename supplies the API identifier; the display name is separate.
    Hover is an expansion capability, not a fifth directory.
    Discovery scans valid Python files without a manual all-list registration.
    Do not introduce a replacement loader.}
\end{frame}

\begin{frame}{Preserve the module interface}
    \begin{itemize}
        \item \code{introspection()}: supported input/output or family-specific setup.
        \item \code{version()}: accurate capabilities, configuration and metadata.
        \item \code{dict_handler(request=...)}: dictionary execution path.
        \item \code{handler(q=...)}: supported legacy JSON execution path.
        \item Return fresh result/error objects; keep request state local.
    \end{itemize}
    \begin{alertblock}{Read the fallback behavior}
        The reviewed dispatcher catches \code{AttributeError} for legacy fallback.
        Do not use that exception as input validation.
    \end{alertblock}
    \note{The reference implements both paths and compares semantic results.
    Avoid mutable shared error dictionaries or metadata modified by a request.
    Importing a module should not contact a provider or change a target system.}
\end{frame}

\begin{frame}{Write down the request and result schema}
    \begin{itemize}
        \item Supported attributes, event envelopes or file inputs?
        \item Settings in \code{config}, \code{params} or import options?
        \item Simplified expansion or \code{format: misp_standard}?
        \item Attributes, objects, an encoded file or action status?
        \item Missing data, empty results and partial failures?
    \end{itemize}
    \begin{block}{Use existing MISP semantics}
        Map the provider's response into valid types, object templates
        and the selected family's contract.
    \end{block}
    \note{Request and response mappings belong in the plan.
    Export modules often return base64 data; actions commonly use params and data,
    but these summaries do not replace reading actual contracts.
    JSON-serializable output can still have the wrong MISP structure.}
\end{frame}

\begin{frame}{Keep provenance and detection intent}
    \begin{itemize}
        \item Preserve the source UUID and original evidence.
        \item Use PyMISP to build and serialize structured objects.
        \item Link derived objects to the correct source.
        \item Choose valid types and object relations.
        \item Do not silently promote contextual data to IDS indicators.
    \end{itemize}
    \begin{block}{Reference exercise}
        Keep the original attribute; add a normalized \code{domain-ip}
        object with a \code{derived-from} reference.
    \end{block}
    \note{Tests inspect real PyMISP output: source UUID, value/comment, references
    and detection flags. The derived attribute is non-IDS in this exercise.
    Validate flags before PyMISP conversion: a nonempty string such as false can
    otherwise become true. The exercise accepts only boolean flags when supplied.
    It is the full normalized name, not a guessed registrable domain.}
\end{frame}

\begin{frame}{Dependencies and configuration are not incidental}
    \begin{itemize}
        \item Use the service's actual Python environment and dependency workflow.
        \item Add necessary dependencies without unrelated lockfile upgrades.
        \item Metadata describing a dependency does not install it.
        \item Distinguish request config, service defaults and module-specific
        environment variables.
        \item Keep credentials and private incidents out of prompts and fixtures.
    \end{itemize}
    \note{Do not bypass artifact or TLS verification to make setup pass.
    Missing optional imports can skip a module; syntax errors can stop startup.
    Neither is the same as MISP enablement or provider authorization.
    Read family-specific configuration and action semantics before implementing.}
\end{frame}

\section{Direct the agent and test behavior}

\begin{frame}{A useful brief needs an explicit contract}
    \begin{enumerate}
        \item Purpose, gap and unique identifier.
        \item Family, supported inputs and expected outputs.
        \item Official provider/schema version and configuration keys.
        \item Side effects and permitted execution boundaries.
        \item Acceptance tests, error cases and required evidence.
    \end{enumerate}
    \begin{block}{Reusable task prompt}
        Start with \code{AGENT_BRIEF.md}.
        It is a prompt to adapt, not an automatically installed agent policy.
    \end{block}
    \note{Ask for a short contract and implementation plan first.
    Continue useful reversible work without repeatedly asking preferences.
    Clarify genuine blockers and keep live submissions or writes within
    the user's authorization.}
\end{frame}

\begin{frame}{Worked exercise: normalize a DNS name locally}
    \begin{itemize}
        \item Input: \code{domain}/\code{hostname}, value and source UUID.
        \item Normalize case, one root dot and standard-library IDNA.
        \item Reject malformed labels, unsupported types, IP literals and bad UUIDs.
        \item Preserve the input; add a referenced MISP-standard object.
        \item No DNS lookup, provider account, model call or runtime agent framework.
    \end{itemize}
    \begin{alertblock}{Training scope}
        This is a contract exercise, not a proposal to add another normalizer upstream.
        A contribution must address a real project gap.
    \end{alertblock}
    \note{The README defines the exact grammar and acceptance cases.
    The exercise uses multi-label names, at most 63 bytes per ASCII label and
    253 bytes overall after IDNA. It does not check registration, reachability,
    public suffixes or maliciousness.}
\end{frame}

\begin{frame}[fragile]{An acceptance request the agent can test}
    \lstinputlisting[language=json]{examples/request.json}
    \small
    Expected: \code{xn--bcher-kva.example} in a \code{domain-ip} object
    linked to the original source UUID.
    \note{The JSON escape represents uppercase U with umlaut.
    Original input value and detection flag stay unchanged in the source Attribute.
    The new object's domain attribute is non-IDS. Check the serialized result,
    not just an agent's printed summary.}
\end{frame}

\begin{frame}{Ask for small, reviewable implementation steps}
    \begin{enumerate}
        \item Metadata and introspection.
        \item Pure parsing/normalization and validation.
        \item MISP result construction with references.
        \item A thin provider boundary, if the task needs one.
        \item Tests and documentation for the actual contract.
    \end{enumerate}
    \begin{block}{Inspect the diff}
        Reject unrelated framework changes, dependency churn,
        import-time networking and invented endpoint behavior.
    \end{block}
    \note{The reference is an answer key for the exercise.
    Compare a real integration with a recent maintained sibling.
    Reuse existing helpers/builders instead of duplicating target-format logic
    or changing the dispatcher to accommodate one module.}
\end{frame}

\begin{frame}{If a module calls a provider, test the boundary}
    \begin{itemize}
        \item Mock the provider client or HTTP boundary where it is used.
        \item Use documented fixtures: valid, empty and malformed responses.
        \item Test authentication errors, rate limits and bounded timeouts.
        \item Preserve references across multiple and partial results.
        \item Verify scope and acknowledgements for operations that write.
    \end{itemize}
    \begin{block}{Keep output construction real}
        Mock external I/O, not the PyMISP objects that assertions need to inspect.
    \end{block}
    \note{The reference has no external I/O to mock. The source audit links
    real upstream DNS tests as an example. Tests should not need real credentials
    by default. A transport success does not establish complete publication or
    activation at the target.}
\end{frame}

\begin{frame}{Make tests challenge the implementation}
    \small
    \begin{tabular}{@{}P{3.5cm}P{6.7cm}@{}}
        \toprule
        Dimension & What can go wrong \\
        \midrule
        Input/schema & Missing fields, unsupported types, Unicode and length boundaries. \\
        MISP semantics & Wrong relations, lost source UUIDs, changed detection flags. \\
        State isolation & Caller or request mutates another result or shared metadata. \\
        Errors & Exceptions hidden as empty success; errors ignored at HTTP 200. \\
        Deployment & Service imports another package or skips the new file. \\
        \bottomrule
    \end{tabular}
    \note{Ask which failure each assertion would catch. Do not weaken assertions
    or delete fixtures just to turn a run green.
    The invalid inputs are subtests; positive cases inspect actual MISP objects.
    Syntax-only tests cannot demonstrate these semantics.}
\end{frame}

\begin{frame}{Report test levels honestly}
    \begin{itemize}
        \item \textbf{Unit}: validation, parsing, serialization and isolation.
        \item \textbf{Service}: discovery and real \code{/query} requests.
        \item \textbf{MISP}: family settings, user access and UI ingestion.
        \item \textbf{Provider/target}: authorization, acknowledgements and activation.
    \end{itemize}
    \begin{block}{Keep evidence scoped}
        Record commands, exit codes and counts.
        Distinguish passed, failed, skipped and unrun checks.
    \end{block}
    \note{The verifier runs eight unit tests by default. Adding --url runs
    three HTTP tests too. It says explicitly when HTTP checks were not requested.
    Unit passes do not imply a tested service, MISP UI or provider integration.}
\end{frame}

\section{Verify and contribute}

\begin{frame}[fragile]{Check what the service actually imports}
    \begin{itemize}
        \item Use an editable package install or explicit custom root.
        \item Install dependencies in the service environment.
        \item Restart the correct process.
    \end{itemize}
\begin{lstlisting}[language=bash]
# From this companion deck's directory:
misp-modules -l 127.0.0.1 -p 6666 \
  --custom "$PWD/examples/custom"
\end{lstlisting}
    The root contains \code{expansion/training_domain_normalize.py}.
    \note{Pass the root, not its expansion directory.
    Upstream modules use the normal package family directory.
    A duplicate custom identifier replaces an installed module, so choose names
    deliberately and use a separate workshop service.}
\end{frame}

\begin{frame}[fragile]{Verify discovery and dispatch}
\begin{lstlisting}[language=bash]
curl --fail --silent --show-error \
  http://127.0.0.1:6666/modules |
  jq '.[] | select(.name == "training_domain_normalize")'

python3 examples/verify_module.py \
  --url http://127.0.0.1:6666 -v
\end{lstlisting}
    \begin{itemize}
        \item Check family, capabilities, input format and config schema.
        \item Verify valid results and invalid-input errors through HTTP.
        \item Inspect \code{/openapi.json}, not just the handler.
    \end{itemize}
    \note{Expected: eight unit plus three service tests, all executed.
    Discovery failures can reflect dependencies, paths, stale installation or
    startup errors. Health alone does not prove the new module loaded.}
\end{frame}

\begin{frame}{HTTP 200 is not the acceptance criterion}
    \begin{itemize}
        \item Inspect \code{error}/\code{errors} and the actual output envelope.
        \item Check references and source fields after serialization.
        \item Decode exports and inspect the target format.
        \item Check action outcomes appropriate to the operation.
        \item Verify MISP enablement separately from service discovery.
    \end{itemize}
    \note{The wildcard example returns a module error over HTTP 200.
    A MISP UI check also needs the relevant family configured, module enabled,
    supported input and organisation access.
    Do not turn a localhost API pass into a claim of production integration.}
\end{frame}

\begin{frame}{Use the project's quality and documentation workflow}
    \begin{itemize}
        \item Read current CI, Poetry groups, extras and native prerequisites.
        \item Run relevant tests and formatting/lint under project configuration.
        \item Broaden checks when changes or failures warrant it.
        \item Supply complete metadata and user documentation.
        \item Keep dependency/lock edits tied to the implementation.
    \end{itemize}
    \sources{The brief supplies command templates; substitute actual paths and family.}
    \note{The reviewed CI installs the test group and all extras, starts a
    service and runs pytest across its matrix. This is broader than the base
    environment used for the reference. Documentation generation also has
    prerequisites; a formatter pass or zero-test run is not upstream CI success.}
\end{frame}

\begin{frame}{What the pull request should explain}
    \begin{enumerate}
        \item Integration gap and resulting behavior.
        \item Input/output examples and configuration, without secrets.
        \item Provider/schema version and side effects.
        \item Actual validation, commands, counts and limitations.
        \item Dependencies and documentation needed by users.
    \end{enumerate}
    \begin{block}{Keep it focused}
        Include the implementation, tests and documentation a maintainer needs.
        Remove incidental edits and generated noise.
    \end{block}
    \note{Describe the final implementation, not an agent conversation or
    abandoned plan. Separate mocked, local-service and live-integration checks.
    An AI-generated assertion that code is ready is not execution evidence.}
\end{frame}

\begin{frame}{Human review is part of the workflow}
    \begin{itemize}
        \item Compare mappings with official API/schema documentation.
        \item Read validation, errors, references and config behavior.
        \item Inspect tests for missing cases, not only their exit status.
        \item Check for unexpected target writes or runtime model calls.
        \item Reconsider a duplicate or unnecessary upstream integration.
    \end{itemize}
    \note{The contributor owns the proposal and maintainers decide on merge.
    Prepare or publish a PR within the user's authorization.
    This exercise does not submit anything to misp-modules upstream; do not
    propose the toy normalizer without a genuine project need.}
\end{frame}

\begin{frame}{Workshop: brief, implement, verify}
    \begin{enumerate}
        \item Adapt \code{AGENT_BRIEF.md} to the exercise.
        \item Ask for a contract and acceptance plan.
        \item Implement in a separate local directory; use the reference as an answer key.
        \item Run unit and service checks and review source/object relationships.
        \item Draft a PR description stating what was and was not tested.
    \end{enumerate}
    \note{Suggested timing: ten minutes for source reading and brief,
    twenty to thirty for implementation/tests, ten for service checks and review.
    Do not create an upstream PR for the toy exercise.
    A real contribution replaces the capstone with an identified missing capability.}
\end{frame}

\begin{frame}{The logic to preserve}
    \begin{enumerate}
        \item Existing loader, family and dispatch contracts.
        \item Types, object templates, source references and detection intent.
        \item Explicit configuration, bounded I/O and useful failures.
        \item Repeatable tests and installation/discovery evidence.
        \item A contribution that a maintainer can review.
    \end{enumerate}
    \note{Use these as a final rubric. An assistant can help at each stage
    without inventing evidence, protocols or an architecture that bypasses
    the project's existing behavior.}
\end{frame}

\begin{frame}{Resources for the next contribution}
    \begin{itemize}
        \item \code{README.md}: setup, exercise specification and source audit.
        \item \code{AGENT_BRIEF.md}: reusable task prompt.
        \item \code{examples/}: reference, request and verifier.
        \item \href{https://github.com/MISP/misp-modules}{MISP/misp-modules} and its current contribution guide.
        \item \href{https://misp.github.io/misp-modules/}{Module and installation documentation}.
        \item The companion \emph{MISP modules} introduction and local API lab.
    \end{itemize}
    \note{Reread the target checkout before the next integration.
    Pinned evidence documents a reviewed snapshot, not a permanent guarantee
    about future APIs, dependencies or CI commands.}
\end{frame}
